\documentclass[10pt]{beamer}
\usetheme{metropolis}
\usepackage{graphicx}
\usepackage{listings}
\usepackage{booktabs}
\usepackage{xcolor}
\definecolor{codebg}{HTML}{F5F5F5}
\definecolor{codegreen}{HTML}{2E7D32}
\definecolor{codegray}{HTML}{757575}
\definecolor{codeblue}{HTML}{1565C0}
\lstdefinestyle{rstyle}{language=R,backgroundcolor=\color{codebg},basicstyle=\ttfamily\scriptsize,
  keywordstyle=\color{codeblue}\bfseries,stringstyle=\color{codegreen},
  commentstyle=\color{codegray}\itshape,breaklines=true,frame=single,
  rulecolor=\color{codegray!30},showstringspaces=false,tabsize=2,
  morekeywords={library,ggplot,aes,geom_bar,geom_col,geom_rect,geom_text,
    coord_polar,scale_fill_manual,scale_fill_brewer,position_fill,
    theme_minimal,theme_void,theme,element_blank,labs,
    treemapify,geom_treemap,geom_treemap_text,waffle}}
\lstdefinestyle{pythonstyle}{language=Python,backgroundcolor=\color{codebg},basicstyle=\ttfamily\scriptsize,
  keywordstyle=\color{codeblue}\bfseries,stringstyle=\color{codegreen},
  commentstyle=\color{codegray}\itshape,breaklines=true,frame=single,
  rulecolor=\color{codegray!30},showstringspaces=false,tabsize=4,
  morekeywords={import,from,as,pd,plt,sns,np,fig,ax,pie,bar,barh,squarify,
    Rectangle,add_patch,Waffle}}
\title{Proportions: Pie, Donut, Treemap \& Waffle}
\subtitle{Workshop 3 --- Module 5}
\author{Prof.Asoc.\ Endri Raco}
\institute{Data Visualization for Data Scientists \\ UNYT Tirana}
\date{2025/26}
\begin{document}

\begin{frame}
  \titlepage
\end{frame}

\begin{frame}{Module Roadmap}
  \begin{enumerate}
    \item The pie chart debate: when pies work and when they fail
    \item Donut charts: single-KPI highlights
    \item Better alternatives: stacked bar and sorted bar with \%
    \item Treemaps: hierarchical part-to-whole
    \item Waffle charts: intuitive unit-based proportions
    \item Proportion chart decision tree
  \end{enumerate}
  \begin{exampleblock}{Learning Outcome}
    You will make an evidence-based decision about whether to use a pie
    chart (rarely) or a better alternative (usually), and produce
    treemaps and waffle charts for part-to-whole displays.
  \end{exampleblock}
\end{frame}

\section{The Pie Chart Debate}

\begin{frame}{Pie vs Bar: Same Data, Different Precision}
  \begin{center}
    \includegraphics[width=0.92\textwidth]{figures_w03m05/pie_vs_bar}
  \end{center}
  \begin{alertblock}{Pro-Tip}
    Angle is the \textbf{worst} perceptual channel (Cleveland \& McGill,
    1984). Humans are $\sim$2$\times$ less accurate judging angles than
    lengths. A sorted horizontal bar with direct \% labels always
    outperforms a pie chart for multi-category comparisons.
  \end{alertblock}
\end{frame}

\begin{frame}{When Pies Work: $\leq$3 Slices}
  \begin{center}
    \includegraphics[width=0.92\textwidth]{figures_w03m05/when_pies_work}
  \end{center}
  \begin{exampleblock}{Data Insight}
    Pies work only when: (a) there are $\leq$3 slices, (b) one slice
    clearly dominates (e.g., 73\% vs 27\%), (c) the goal is to show
    ``more than half'' rather than precise comparison. For 4+ slices,
    the angle discrimination fails completely.
  \end{exampleblock}
\end{frame}

\section{Donut Charts}

\begin{frame}{Donut: Use Only for Single-KPI Highlights}
  \begin{center}
    \includegraphics[width=0.82\textwidth]{figures_w03m05/donut}
  \end{center}
\end{frame}

\begin{frame}[fragile]{Donut in R vs Python}
  \begin{columns}[T]
    \begin{column}{0.48\textwidth}
      \textbf{R (ggplot2 + coord\_polar)}
\begin{lstlisting}[style=rstyle]
# Single-KPI donut
df <- tibble(
  cat = c("Free", "Paid"),
  n = c(73, 27))

ggplot(df, aes(x = 2, y = n,
               fill = cat)) +
  geom_col(width = 1,
    color = "white", linewidth = 2) +
  coord_polar(theta = "y") +
  xlim(0.5, 2.5) +   # hole
  scale_fill_manual(values =
    c(Free = "#1565C0",
      Paid = "#EEEEEE")) +
  annotate("text", x = 0.5, y = 0,
    label = "73%", size = 10,
    fontface = "bold",
    color = "#1565C0") +
  theme_void() +
  theme(legend.position = "none")
\end{lstlisting}
    \end{column}
    \begin{column}{0.48\textwidth}
      \textbf{Python}
\begin{lstlisting}[style=pythonstyle]
fig, ax = plt.subplots()

wedges, _ = ax.pie(
    [73, 27],
    colors=["#1565C0", "#EEEEEE"],
    startangle=90,
    wedgeprops={
        "width": 0.35,     # ring
        "edgecolor": "white",
        "linewidth": 2})

# Center text = the KPI
ax.text(0, 0, "73%",
    ha="center", va="center",
    fontsize=28, fontweight="bold",
    color="#1565C0")
ax.text(0, -0.12, "Free Apps",
    ha="center", va="center",
    fontsize=10, color="#555")
\end{lstlisting}
    \end{column}
  \end{columns}
\end{frame}

\section{Better Alternatives}

\begin{frame}{Stacked Bar and Sorted Bar with \%}
  \begin{center}
    \includegraphics[width=0.92\textwidth]{figures_w03m05/stacked_alternatives}
  \end{center}
\end{frame}

\begin{frame}[fragile]{100\% Stacked Bar in Code}
  \begin{columns}[T]
    \begin{column}{0.48\textwidth}
      \textbf{R}
\begin{lstlisting}[style=rstyle]
# 100% stacked (fill)
ggplot(apps, aes(x = "Apps",
    fill = Category)) +
  geom_bar(position = "fill",
    width = 0.6) +
  geom_text(aes(
    label = paste0(
      round(..count../sum(..count..)*
        100), "%")),
    stat = "count",
    position = position_fill(
      vjust = 0.5),
    size = 3, color = "white",
    fontface = "bold") +
  scale_y_continuous(
    labels = scales::percent) +
  coord_flip() +
  theme_minimal()
\end{lstlisting}
    \end{column}
    \begin{column}{0.48\textwidth}
      \textbf{Python}
\begin{lstlisting}[style=pythonstyle]
# Manual 100% stacked bar
left = 0
for cat, val, color in zip(
    cats, vals, colors):
    pct = val / sum(vals) * 100
    ax.barh(["Apps"], pct,
        left=left, color=color,
        height=0.5,
        edgecolor="white", lw=1.5,
        label=cat)
    if pct > 5:
        ax.text(left + pct/2, 0,
            f"{pct:.0f}%",
            ha="center", va="center",
            fontsize=7,
            fontweight="bold",
            color="white")
    left += pct
ax.set_xlim(0, 100)
ax.legend(fontsize=6, ncol=3)
\end{lstlisting}
    \end{column}
  \end{columns}
\end{frame}

\section{Treemaps}

\begin{frame}{Treemap: Area Encodes Magnitude}
  \begin{center}
    \includegraphics[width=0.78\textwidth]{figures_w03m05/treemap}
  \end{center}
  \begin{exampleblock}{Data Insight}
    Treemaps encode part-to-whole via area, which is more accurate than
    angle (pie) but less accurate than length (bar). Their advantage is
    showing \textbf{hierarchical} structure: category $\rightarrow$
    sub-category $\rightarrow$ individual. Use for hierarchical data with
    10--50 leaf nodes.
  \end{exampleblock}
\end{frame}

\begin{frame}[fragile]{Treemap in R vs Python}
  \begin{columns}[T]
    \begin{column}{0.48\textwidth}
      \textbf{R (treemapify)}
\begin{lstlisting}[style=rstyle]
library(treemapify)

ggplot(df, aes(
    area = n,
    fill = Category,
    label = paste(Category,
      scales::comma(n),
      sep = "\n"))) +
  geom_treemap(color = "white",
    size = 2) +
  geom_treemap_text(
    fontface = "bold",
    color = "white",
    place = "centre") +
  scale_fill_brewer(
    palette = "Set2") +
  theme(legend.position = "none") +
  labs(title = "App Categories")
\end{lstlisting}
    \end{column}
    \begin{column}{0.48\textwidth}
      \textbf{Python (squarify)}
\begin{lstlisting}[style=pythonstyle]
import squarify

fig, ax = plt.subplots(
    figsize=(8, 5))

squarify.plot(
    sizes=vals,
    label=[f"{c}\n{v:,}"
        for c, v in
        zip(cats, vals)],
    color=colors,
    alpha=0.7,
    edgecolor="white",
    linewidth=2,
    text_kwargs={
        "fontsize": 8,
        "fontweight": "bold",
        "color": "white"},
    ax=ax)

ax.axis("off")
ax.set_title("Treemap")
\end{lstlisting}
    \end{column}
  \end{columns}
\end{frame}

\section{Waffle Charts}

\begin{frame}{Waffle: 1 Square = 1\%}
  \begin{center}
    \includegraphics[width=0.72\textwidth]{figures_w03m05/waffle}
  \end{center}
  \begin{alertblock}{Pro-Tip}
    Waffle charts use a 10$\times$10 grid where each square represents 1\%.
    They are more intuitive than pies because counting squares is easier
    than judging angles. In R: \texttt{waffle} package. In Python:
    \texttt{pywaffle} or manual \texttt{plt.Rectangle()} loop.
    Limit to $\leq$6 categories for readability.
  \end{alertblock}
\end{frame}

\begin{frame}[fragile]{Waffle in R vs Python}
  \begin{columns}[T]
    \begin{column}{0.48\textwidth}
      \textbf{R (waffle)}
\begin{lstlisting}[style=rstyle]
# install.packages("waffle")
library(waffle)

waffle(
  c(FAMILY = 38, GAME = 20,
    TOOLS = 18, OTHER = 24),
  rows = 10,
  colors = c("#1565C0","#E53935",
    "#2E7D32","#BBBBBB"),
  title = "App Categories (%)",
  xlab = "1 square = 1%"
)
\end{lstlisting}
    \end{column}
    \begin{column}{0.48\textwidth}
      \textbf{Python (manual)}
\begin{lstlisting}[style=pythonstyle]
fig, ax = plt.subplots()
# Build 10x10 grid
grid = np.zeros((10,10), dtype=int)
idx = 0; cat_idx = 0
for i in range(10):
    for j in range(10):
        grid[9-i, j] = cat_idx
        idx += 1
        if (idx >= sum(pcts[:cat_idx+1])
            and cat_idx < len(pcts)-1):
            cat_idx += 1

# Draw squares
for i in range(10):
    for j in range(10):
        ax.add_patch(plt.Rectangle(
            (j, i), 0.9, 0.9,
            fc=colors[grid[i,j]],
            ec="white", lw=1.5))
ax.set_aspect("equal")
ax.axis("off")
\end{lstlisting}
    \end{column}
  \end{columns}
\end{frame}

\section{Decision Framework}

\begin{frame}{Proportion Chart Decision Tree}
  \begin{center}
    \includegraphics[width=0.88\textwidth]{figures_w03m05/decision_tree}
  \end{center}
\end{frame}

\begin{frame}{Proportion Chart Comparison}
  \begin{center}
    \includegraphics[width=0.92\textwidth]{figures_w03m05/comparison_card}
  \end{center}
\end{frame}

\section{Complete Examples}

\begin{frame}[fragile]{R: Single-KPI Donut}
  \begin{columns}[T]
    \begin{column}{0.50\textwidth}
\begin{lstlisting}[style=rstyle]
# The only acceptable pie/donut:
# 2 slices + center KPI
df <- tibble(
  cat = c("Free", "Other"),
  n = c(73, 27))

ggplot(df, aes(x = 2, y = n,
               fill = cat)) +
  geom_col(width = 1,
    color = "white", lw = 2) +
  coord_polar(theta = "y") +
  xlim(0.5, 2.5) +
  scale_fill_manual(values =
    c(Free = "#1565C0",
      Other = "#EEEEEE")) +
  annotate("text", x = 0.5, y = 0,
    label = "73%", size = 12,
    fontface = "bold",
    color = "#1565C0") +
  annotate("text", x = 0.5, y = 5,
    label = "Free Apps",
    size = 4, color = "#555") +
  theme_void() +
  theme(legend.position = "none")
\end{lstlisting}
    \end{column}
    \begin{column}{0.45\textwidth}
      \includegraphics[width=\textwidth]{figures_w03m05/r_output}
    \end{column}
  \end{columns}
\end{frame}

\begin{frame}[fragile]{Python: Waffle Chart}
  \begin{columns}[T]
    \begin{column}{0.50\textwidth}
\begin{lstlisting}[style=pythonstyle]
# Each square = 1%
# 10x10 grid = 100 squares

pcts = [38, 20, 18, 24]  # rounded
cats = ["FAMILY","GAME",
        "TOOLS","OTHER"]
colors = ["#1565C0","#E53935",
          "#2E7D32","#BBBBBB"]

fig, ax = plt.subplots(
    figsize=(6, 5))
# ... grid construction ...
# ... Rectangle loop ...

# Legend
for k, (cat, pct, c) in enumerate(
    zip(cats, pcts, colors)):
    ax.text(10.8, 8.5 - k*1.2,
        f"  {cat} ({pct}%)",
        fontsize=8, color=c,
        fontweight="bold")

ax.set_title(
    "1 square = 1%",
    fontweight="bold")
\end{lstlisting}
    \end{column}
    \begin{column}{0.45\textwidth}
      \includegraphics[width=\textwidth]{figures_w03m05/py_output}
    \end{column}
  \end{columns}
\end{frame}

\section{Synthesis}

\begin{frame}{Key Takeaways}
  \begin{enumerate}
    \item \textbf{Pie charts fail} for 4+ slices because angle is the
          worst perceptual channel. Use only for 2--3 slices with clear
          dominance.
    \item \textbf{Donut charts} are pies with a hole --- same problems,
          but acceptable for a single-KPI highlight with center text.
    \item \textbf{100\% stacked bar} or \textbf{sorted bar with \%}
          almost always outperform pies for multi-category proportions.
    \item \textbf{Treemaps} encode part-to-whole via area; best for
          hierarchical data with 10--50 nodes.
    \item \textbf{Waffle charts} (1 square = 1\%) are more intuitive
          than pies for general audiences.
    \item Default to \textbf{sorted horizontal bar with \%} --- it is
          the safest proportion chart for any audience.
  \end{enumerate}
\end{frame}

\begin{frame}{Essential Readings}
  \begin{itemize}
    \item Cleveland, W.~S. \& McGill, R. (1984). ``Graphical Perception.''
          \emph{JASA}, 79(387). (Angle $<$ length for proportion.)
    \item Spence, I. \& Lewandowsky, S. (1991). ``Displaying proportions
          and percentages.'' \emph{Applied Cognitive Psychology}.
    \item Schwabish, J. (2021). \emph{Better Data Visualizations},
          Chapter 6 (Part-to-Whole).
    \item Wilke, C.~O. (2019). \emph{Fundamentals of Data Visualization},
          Chapter 10 (Proportions).
  \end{itemize}
  \vspace{6pt}
  \textbf{Next Module}: Rankings \& Part-to-Whole (W03-M06)
\end{frame}

\end{document}
